% Options for packages loaded elsewhere
\PassOptionsToPackage{unicode}{hyperref}
\PassOptionsToPackage{hyphens}{url}
\documentclass[
  ignorenonframetext,
  aspectratio=169,
]{beamer}
\newif\ifbibliography
\usepackage{pgfpages}
\setbeamertemplate{caption}[numbered]
\setbeamertemplate{caption label separator}{: }
\setbeamercolor{caption name}{fg=normal text.fg}
\beamertemplatenavigationsymbolsempty
% remove section numbering
\setbeamertemplate{part page}{
  \centering
  \begin{beamercolorbox}[sep=16pt,center]{part title}
    \usebeamerfont{part title}\insertpart\par
  \end{beamercolorbox}
}
\setbeamertemplate{section page}{
  \centering
  \begin{beamercolorbox}[sep=12pt,center]{section title}
    \usebeamerfont{section title}\insertsection\par
  \end{beamercolorbox}
}
\setbeamertemplate{subsection page}{
  \centering
  \begin{beamercolorbox}[sep=8pt,center]{subsection title}
    \usebeamerfont{subsection title}\insertsubsection\par
  \end{beamercolorbox}
}
% Prevent slide breaks in the middle of a paragraph
\widowpenalties 1 10000
\raggedbottom
\AtBeginPart{
  \frame{\partpage}
}
\AtBeginSection{
  \ifbibliography
  \else
    \frame{\sectionpage}
  \fi
}
\AtBeginSubsection{
  \frame{\subsectionpage}
}
\usepackage{iftex}
\ifPDFTeX
  \usepackage[T1]{fontenc}
  \usepackage[utf8]{inputenc}
  \usepackage{textcomp} % provide euro and other symbols
\else % if luatex or xetex
  \usepackage{unicode-math} % this also loads fontspec
  \defaultfontfeatures{Scale=MatchLowercase}
  \defaultfontfeatures[\rmfamily]{Ligatures=TeX,Scale=1}
\fi
\usepackage{lmodern}
\ifPDFTeX\else
  % xetex/luatex font selection
\fi
% Use upquote if available, for straight quotes in verbatim environments
\IfFileExists{upquote.sty}{\usepackage{upquote}}{}
\IfFileExists{microtype.sty}{% use microtype if available
  \usepackage[]{microtype}
  \UseMicrotypeSet[protrusion]{basicmath} % disable protrusion for tt fonts
}{}
\makeatletter
\@ifundefined{KOMAClassName}{% if non-KOMA class
  \IfFileExists{parskip.sty}{%
    \usepackage{parskip}
  }{% else
    \setlength{\parindent}{0pt}
    \setlength{\parskip}{6pt plus 2pt minus 1pt}}
}{% if KOMA class
  \KOMAoptions{parskip=half}}
\makeatother
\usepackage{longtable,booktabs,array}
\usepackage{caption}
\captionsetup[table]{skip=6pt}
\newcounter{none} % for unnumbered tables
\usepackage{calc} % for calculating minipage widths
\usepackage{caption}
% Make caption package work with longtable
\makeatletter
\def\fnum@table{\tablename~\thetable}
\makeatother
\setlength{\emergencystretch}{3em} % prevent overfull lines
\providecommand{\tightlist}{%
  \setlength{\itemsep}{0pt}\setlength{\parskip}{0pt}}
% Auto-generated by infrastructure/rendering/_pdf_latex_helpers.py::extract_math_font_preamble.
% Minimal header passed to Pandoc via -H so Beamer slides pick up the same math font
% as the combined-PDF path. Adjust the manuscript's preamble.md to change the font globally.
\usepackage{unicode-math}
\setmathfont{latinmodern-math.otf}

% Manuscript macro/package fallbacks (newcommand -> providecommand).
% Core mathematics
\usepackage{amsmath}
\usepackage{amssymb}
% Document layout
% Tables
\usepackage{booktabs}
% Caption styling (booktabs already loaded above). Small caption font with a
% bold label ("Figure 1", "Table 2") gives the math/figure-heavy layout a
% consistent, journal-like caption rhythm without fighting pandoc-crossref —
% pandoc emits standard \caption{}, and caption/captionsetup only restyle it.
% ── Theorem environments ─────────────────────────────────────────────────
% amsthm is loaded above. The FedGVI<->Friston bridge is stated as a small
% number of theorems/definitions/lemmas/propositions/corollaries; sharing one
% counter (definition/lemma/proposition/corollary all step the `theorem`
% counter) gives a single monotone numbering 1, 2, 3, ... across all kinds, so
% authors can reference them as Theorem 1, Definition 2, Corollary 3 without a
% per-kind counter drifting out of order. Plain style for theorem-like results,
% definition style (upright body) for definitions.
% (algorithm2e intentionally not loaded: this manuscript states results as
% numbered amsthm Theorems/Definitions, not algorithm2e pseudocode.)
% Code listings
\usepackage{listings}
% Typography and formatting
% (siunitx intentionally not loaded: no \SI/\num units appear in this manuscript.)
% Cross-references and citations
% (cleveref and titlesec intentionally not loaded: unavailable in this TeX tree
% and not required. Cross-references resolve through pandoc-crossref's own
% \ref-style names — no \cref appears — and section heading styling falls back
% to the pandoc/LaTeX defaults.)
% ── TeX Live Basic-compatible mono font for code listings ────────────
% Latin Modern Mono is bundled with TeX Live Basic and keeps the manuscript
% renderable on minimal TeX installations. Mathematical Unicode coverage is
% provided separately by Latin Modern Math below, so the code font does not
% need a private JuliaMono installation.
%
% Use the TeX font filename rather than the system-family name: the minimal
% TeX Live fontconfig database may not register the OpenType family even though
% kpsewhich can resolve the bundled files.
% Math font for unicode-math: Latin Modern Math (TeX Live) has full BMP
% coverage including U+2223 (\mid), U+226A/226B (\ll/\gg), and the Greek/
% blackboard letters used in equations. Without an explicit \setmathfont,
% unicode-math falls back to lmroman text font which lacks several glyphs
% and emits "Missing character" warnings on every \mid in math mode.

% Slide layout defaults for warning-clean scientific decks.
\usepackage{etoolbox}
\IfFileExists{xurl.sty}{\usepackage{xurl}}{}
\IfFileExists{seqsplit.sty}{\usepackage{seqsplit}}{\newcommand{\seqsplit}[1]{#1}}
\protected\def\breakseq#1{\seqsplit{#1}}
\protected\def\breaktt#1{\begingroup\ttfamily\seqsplit{#1}\endgroup}
\setlength{\emergencystretch}{6em}
\tolerance=5000
\hbadness=10000
\hfuzz=1pt
\setlength{\tabcolsep}{2pt}
\AtBeginEnvironment{longtable}{\tiny\renewcommand{\arraystretch}{0.86}\setlength{\tabcolsep}{1pt}}
\AtBeginEnvironment{tabular}{\tiny\renewcommand{\arraystretch}{0.86}\setlength{\tabcolsep}{1pt}}
\AtBeginEnvironment{equation}{\tiny}
\AtBeginEnvironment{equation*}{\tiny}
\AtBeginEnvironment{align}{\tiny}
\AtBeginEnvironment{align*}{\tiny}
\AtBeginEnvironment{itemize}{\footnotesize}
\AtBeginEnvironment{enumerate}{\footnotesize}
\AtBeginEnvironment{description}{\footnotesize}
\setbeamerfont{caption}{size=\tiny}
\setbeamerfont{caption name}{size=\tiny}
\setbeamerfont{normal text}{size=\small}
\setbeamerfont{frametitle}{size=\small}
\setbeamerfont{section title}{size=\footnotesize}
\setbeamerfont{subsection title}{size=\footnotesize}
\setbeamertemplate{section page}{%
  \centering
  \begin{beamercolorbox}[sep=12pt,center,wd=\paperwidth]{section title}
    \parbox{0.86\paperwidth}{\centering\usebeamerfont{section title}\insertsection\par}
  \end{beamercolorbox}
}
\setbeamertemplate{subsection page}{%
  \centering
  \begin{beamercolorbox}[sep=8pt,center,wd=\paperwidth]{subsection title}
    \parbox{0.86\paperwidth}{\centering\usebeamerfont{subsection title}\insertsubsection\par}
  \end{beamercolorbox}
}
\setlength{\abovecaptionskip}{2pt}
\setlength{\belowcaptionskip}{0pt}

% Natbib and cross-reference fallbacks — slides don't load natbib
% or cleveref, but combined-PDF manuscript prose may emit these
% commands. The fallback renders citations as a bracketed key list
% and cross-references as detokenized labels so slides don't fail on
% undefined control sequences or raw underscores. \providecommand is
% a no-op if packages load later.
\providecommand{\citep}[1]{[#1]}
\providecommand{\citet}[1]{#1}
\providecommand{\citealp}[1]{#1}
\providecommand{\citeauthor}[1]{#1}
\providecommand{\citeyear}[1]{#1}
\providecommand{\cref}[1]{\texttt{\detokenize{#1}}}
\providecommand{\Cref}[1]{\texttt{\detokenize{#1}}}
\providecommand{\crefrange}[2]{\texttt{\detokenize{#1}}--\texttt{\detokenize{#2}}}
\providecommand{\Crefrange}[2]{\texttt{\detokenize{#1}}--\texttt{\detokenize{#2}}}
\providecommand{\paragraph}[1]{\textbf{#1}\ }

% Opt-in accessible presentation profile.
\setbeamerfont{normal text}{size*={20pt}{24pt}}
\setbeamerfont{frametitle}{size*={28pt}{32pt}}
\setbeamerfont{section title}{size*={28pt}{32pt}}
\setbeamerfont{subsection title}{size*={28pt}{32pt}}
\setbeamerfont{caption}{size*={16pt}{19pt}}
\setbeamerfont{caption name}{size*={16pt}{19pt}}
\setbeamerfont{itemize/enumerate body}{size*={20pt}{24pt}}
\setbeamerfont{itemize/enumerate subbody}{size*={20pt}{24pt}}
\setbeamerfont{itemize/enumerate subsubbody}{size*={20pt}{24pt}}
\setbeamerfont{itemize item}{size*={20pt}{24pt}}
\setbeamerfont{itemize subitem}{size*={20pt}{24pt}}
\setbeamerfont{itemize subsubitem}{size*={20pt}{24pt}}
\setbeamerfont{description body}{size*={20pt}{24pt}}
\setbeamerfont{description item}{size*={20pt}{24pt}}
\setbeamerfont{quote}{size*={20pt}{24pt}}
\setbeamerfont{quotation}{size*={20pt}{24pt}}
\AtBeginEnvironment{longtable}{\fontsize{20pt}{24pt}\selectfont}
\AtBeginEnvironment{tabular}{\fontsize{20pt}{24pt}\selectfont}
\AtBeginEnvironment{equation}{\fontsize{20pt}{24pt}\selectfont}
\AtBeginEnvironment{equation*}{\fontsize{20pt}{24pt}\selectfont}
\AtBeginEnvironment{align}{\fontsize{20pt}{24pt}\selectfont}
\AtBeginEnvironment{align*}{\fontsize{20pt}{24pt}\selectfont}
\AtBeginEnvironment{itemize}{\fontsize{20pt}{24pt}\selectfont}
\AtBeginEnvironment{enumerate}{\fontsize{20pt}{24pt}\selectfont}
\AtBeginEnvironment{description}{\fontsize{20pt}{24pt}\selectfont}
\AtBeginDocument{\usebeamerfont{normal text}}
\newlength{\AFSlideTableWidth}
% The fixed footer reservation is calibrated with the semantic seven-line regular-body budget.
\setbeamertemplate{footline}{%
  \leavevmode\vbox{%
    \hbox{%
      \begin{beamercolorbox}[wd=\paperwidth,ht=16pt,dp=3pt,center]{author in head/foot}%
        \fontsize{16pt}{19pt}\selectfont Untagged PDF derivative \textbar\ \href{../web/index.html}{HTML reader}%
      \end{beamercolorbox}%
    }%
    \vskip6pt%
  }%
}
% Accessible footer reserved height: 25pt.

\newtheorem{proposition}{Proposition}
\newtheorem{hypothesis}{Hypothesis}
\newtheorem{remark}[theorem]{Remark}
\newtheorem{axiom}{Axiom}
\newtheorem{property}{Property}
\usepackage{bookmark}
\IfFileExists{xurl.sty}{\usepackage{xurl}}{} % add URL line breaks if available
\urlstyle{same}
% fallback for those not using the hyperref driver hyperxmp:
\makeatletter
\@ifundefined{xmpquote}{\newcommand{\xmpquote}[1]{#1}}{}
\makeatother
\hypersetup{
  hidelinks,
  pdfcreator={LaTeX via pandoc}}

\author{\texorpdfstring{}{}}
\date{}

\begin{document}

\begin{frame}{Supplement: parameter-sensitivity methods}
\protect\phantomsection\label{sec:supp-sensitivity}
This supplement documents how the sensitivity grid of
sec.~14.13 is generated and --- importantly for
a reader trying to reconcile numbers across studies ---
\end{frame}

\begin{frame}{Supplement\ldots{} (part 2)}
where its two sweeps use \emph{different} seeding and trial budgets. It
answers two questions the main section leaves implicit: exactly which
seed drives each cell, so every cell is deterministically reproducible
in isolation;
\end{frame}

\begin{frame}{Supplement\ldots{} (part 3)}
and why the cross-study summary's sensitivity row is not directly
comparable, at matched trial counts, to the standalone heatmap.
\end{frame}

\begin{frame}[fragile]{Experimental protocol for grid sensitivity}
\protect\phantomsection\label{experimental-protocol-for-grid-sensitivity}
The sensitivity sweep uses \texttt{run\_belief\_sharing\_sensitivity}
and \texttt{run\_hierarchical\_sensitivity} from the
\texttt{fedference.experiments} module.
\end{frame}

\begin{frame}{Experimental protocol for\ldots{} (part 2)}
Each function accepts a tuple of acuity values and a tuple of colony
sizes. In the \textbf{belief-sharing} sweep every (acuity, colony-size)
cell averages \(n_{\text{trials}}\) independent trials, each seeded via
a deterministic formula:
\end{frame}

\begin{frame}{Experimental protocol for\ldots{} (part 3)}
\begin{equation}\tag{40}\protect\phantomsection\label{eq:sensitivity-seed}{
\text{seed}_{\text{cell}} = \text{seed}_{\text{base}} + i \cdot 10^5 + j \cdot 10^3 + t
}\end{equation}
\end{frame}

\begin{frame}{Experimental protocol for\ldots{} (part 4)}
where \(i\) indexes acuity, \(j\) indexes colony size, and \(t\) indexes
the trial within a cell. The deterministic rule
eq.~40 makes every grid cell and trial
reproducible in isolation from the base seed. For the belief-sharing
sweep, no two cells share a trial seed,
\end{frame}

\begin{frame}[fragile]{Experimental protocol for\ldots{} (part 5)}
so the design introduces no seed reuse between cells.

Re-running with the same \texttt{seed\_base} is bit-identical. Changing
\texttt{seed\_base} produces an independent replicate for robustness
checking.
\end{frame}

\begin{frame}[fragile]{Experimental protocol for\ldots{} (part 6)}
The \textbf{hierarchical} sweep uses a simpler protocol: every cell
calls \texttt{run\_hierarchical\_world} once with the same base seed
(its internal trials are seeded by that run), so hierarchical cells
share the base seed rather than the per-cell formula above.
\end{frame}

\begin{frame}{Grid parameters for acuity and colony size}
\protect\phantomsection\label{grid-parameters-for-acuity-and-colony-size}
{\def\LTcaptype{none} % do not increment counter
% template-accessible-table-inset
\begingroup
\setlength{\AFSlideTableWidth}{\textwidth}%
\setlength{\LTleft}{\fill}%
\setlength{\LTright}{\fill}%
\begin{longtable}[]{@{}
  >{\raggedright\arraybackslash}p{(\AFSlideTableWidth - 2\tabcolsep) * \real{0.4615}}
  >{\raggedright\arraybackslash}p{(\AFSlideTableWidth - 2\tabcolsep) * \real{0.5385}}@{}}
\toprule\noalign{}
\begin{minipage}[b]{\linewidth}\raggedright
Parameter
\end{minipage} & \begin{minipage}[b]{\linewidth}\raggedright
Values
\end{minipage} \\
\midrule\noalign{}
\endhead
Sensor acuity \(\kappa\) & \{0.40, 0.55, 0.70, 0.85, 0.95\} \\
Colony size \(n\) & \{2, 4, 6, 8, 10\} \\
\bottomrule\noalign{}
\end{longtable}
\endgroup
}
\end{frame}

\begin{frame}[fragile]{Grid parameters for acuity\ldots{} (part 2)}
The 5\texorpdfstring{\ensuremath{\times}}{x}5 = 25 cells per system are run with \texttt{seed\_base\ =\ 0} by
default; \texttt{generate\_sensitivity\_heatmap} accepts a \texttt{seed}
argument to override this.
\end{frame}

\begin{frame}[fragile]{Belief-sharing condition in the sensitivity grid}
\protect\phantomsection\label{belief-sharing-condition-in-the-sensitivity-grid}
Each trial follows four labeled operations. \textbf{Draw:} sample one
true state and one noisy observation per agent using the
\texttt{run\_belief\_sharing} protocol.
\end{frame}

\begin{frame}[fragile]{Belief-sharing condition in\ldots{} (part 2)}
\textbf{Compare:} run one belief-sharing round with
\texttt{communicate=True} and the matched isolated condition with
\texttt{communicate=False}. \textbf{Record:} retain
\texttt{mean\_accuracy} for both conditions.
\end{frame}

\begin{frame}[fragile]{Belief-sharing condition in\ldots{} (part 3)}
\textbf{Reduce:} define the cell value as the mean
communicating-minus-isolated gap across \texttt{n\_trials}.
\end{frame}

\begin{frame}[fragile]{Hierarchical POMDP condition in the sensitivity
grid}
\protect\phantomsection\label{hierarchical-pomdp-condition-in-the-sensitivity-grid}
Each cell in the hierarchical sweep calls
\texttt{run\_hierarchical\_world} once with the cell's acuity and colony
size, passing the constant base seed (not the per-cell formula, which
applies only to the belief-sharing sweep).
\end{frame}

\begin{frame}[fragile]{Hierarchical POMDP\ldots{} (part 2)}
The returned \texttt{location\_accuracy\_gap} (hierarchical minus flat)
becomes the cell value.
\end{frame}

\begin{frame}[fragile]{Figure rendering for sensitivity summaries}
\protect\phantomsection\label{figure-rendering-for-sensitivity-summaries}
\texttt{generate\_sensitivity\_heatmap} assembles the two grids into
vertically stacked matplotlib \texttt{imshow} panels with a
color-vision-deficiency-safer blue--neutral--brown signed palette and
symmetric bounds at \(\pm\max(|\text{gap}|)\),
\end{frame}

\begin{frame}[fragile]{Figure rendering for\ldots{} (part 2)}
per-cell numeric annotations, and a per-panel colorbar labeled
``Accuracy difference''. The figure is written as
\texttt{sensitivity\_heatmap.png} under \texttt{output/figures/}.
Negative gaps map toward blue, positive gaps toward brown,
\end{frame}

\begin{frame}{Figure rendering for\ldots{} (part 3)}
and zero to the neutral midpoint; exact signed cell labels duplicate
that encoding. Hatching marks only the declared
\(|\mathrm{gap}| \le 0.05\) display band. It is not an unreliability
flag, confidence interval, significance test,

or proof of zero effect.
\end{frame}

\begin{frame}[fragile]{Cross-study summary construction}
\protect\phantomsection\label{cross-study-summary-construction}
\texttt{generate\_cross\_study\_summary} performs a separate harmonized
128-seed (\(n_{\text{seeds}} = 128\)) rerun over Studies 1--9 and
reports the mean with a 95\% seed-level percentile-bootstrap interval
for each headline signed contrast or estimand. These intervals belong to
that rerun,
\end{frame}

\begin{frame}{Cross-study summary\ldots{} (part 2)}
including for rows whose primary figure is a deterministic
single-posterior or fixed-configuration diagnostic; they are not
intervals on the primary deterministic display. The definitions are:
\end{frame}

\begin{frame}{Cross-study summary\ldots{} (part 3)}
The robustness row uses 40 matched trials per seed and rate; the
trial-level observations are reduced within seed before the cross-study
summary is formed. Within each seed, the row then takes the
display-selected maximum non-reference server-preset contrast at the
declared worst rate.
\end{frame}

\begin{frame}{Cross-study summary\ldots{} (part 4)}
This preserves the seed as the independent Monte Carlo unit but makes
Study 4 a within-run display-selected summary, not a preselected method
or inferential winner.
\end{frame}

\begin{frame}[fragile]{Cross-study summary\ldots{} (part 5)}
The Study 8 row below uses 3 trials per cell --- smaller than the
full-resolution 20-trial \texttt{Trials\ per\ cell} grid documented
above for the standalone sensitivity heatmap figure, a deliberate
runtime budget for the per-seed cross-study loop rather than an
oversight ---
\end{frame}

\begin{frame}{Cross-study summary\ldots{} (part 6)}
so the two are not directly comparable at matched trial counts.
\end{frame}

\begin{frame}{Cross-study summary\ldots{} (part 7)}
{\def\LTcaptype{none} % do not increment counter
% template-accessible-table-inset
\begingroup
\setlength{\AFSlideTableWidth}{\textwidth}%
\setlength{\LTleft}{\fill}%
\setlength{\LTright}{\fill}%
\begin{longtable}[]{@{}
  >{\raggedright\arraybackslash}p{(\AFSlideTableWidth - 2\tabcolsep) * \real{0.4667}}
  >{\raggedright\arraybackslash}p{(\AFSlideTableWidth - 2\tabcolsep) * \real{0.5333}}@{}}
\toprule\noalign{}
\begin{minipage}[b]{\linewidth}\raggedright
Study
\end{minipage} & \begin{minipage}[b]{\linewidth}\raggedright
Metric
\end{minipage} \\
\midrule\noalign{}
\endhead
1 --- Belief sharing & Accuracy contrast: communicating - isolated \\
2 --- Language acquisition & KL reduction: initial - final \\
\bottomrule\noalign{}
\end{longtable}
\endgroup
}
\end{frame}

\begin{frame}[fragile]{Cross-study summary\ldots{} (part 8)}
Bootstrap CIs use 5000 resamples. Their default \texttt{n\_boot} is
defined by \texttt{bootstrap\_ci} in the \texttt{fedference.statistics}
module.
\end{frame}

\end{document}
